% 第 32 章。数字与教学模型：checks/check_ch32.py；图：figsrc/fig_ch32.py。
% 史料：generaltheory_facts；英国的失业、产出与贴现率存 boe_millennium_unemp_1919_1938_monthly、boe_millennium_gb_1919_1938 与 boe_millennium_bank_rate_1929_1939（取数脚本 data/fetch/boe_millennium_gb32.py）。
\chapter{凯恩斯《通论》}\label{ch:generaltheory}

\begin{chapterguide}
\guidequestion{如果人人都想多储蓄，总储蓄为什么可能反而减少？}
\begin{guidegoals}
\item 用凯恩斯交叉推导有效需求与乘数，并计算节俭悖论；
\item 解释流动性偏好理论，以及它对数量论与可贷资金理论的挑战；
\item 说明“动物精神”与预期在投资中的作用；
\item 区分凯恩斯本人的论证与后来的教科书版本。
\end{guidegoals}
\guideprereq{\chref{ch:say}\enspace\nameref{ch:say}；\chref{ch:depressionus}\enspace\nameref{ch:depressionus}}
\end{chapterguide}

1931 年 1 月，凯恩斯（John Maynard Keynes，1883—1946）在电台上作了一次广播讲话，题目是《储蓄与花费》\cite{keynes1931b,loc32}。讲稿登在 1 月 14 日的《听众》周刊上\cite{keynes1931b}。他估计，英国、德国与美国这三个主要的工业国，大概有 1,200 万产业工人没有工作\cite{keynes1931b}。按政府的行政记录，英国那个月的失业率是 21.1\%\cite{boe2017}。许多关心国家的人以为，他们能做的最有用的事，是比平常多储蓄。凯恩斯说：“在某些情况下，这些都完全对；可是在眼下的情况下，很不幸，它完全错了。”\cite{keynes1931b}
% 据: ss_toc, ss_listener, ss_broadcast, loc_keynes, ss_12m, ss_savemore, ss_wrong

他接着说：“我能做的最好的估计是，你每储蓄五先令，就让一个人失业一天。”\cite{keynes1931b}先令是英国的货币单位，一镑合 20 先令，一先令合 12 便士。五先令就是四分之一镑。照他的说法，一个人失了业，花钱的能力跟着减少，本来要为他生产东西的人又会失业\cite{keynes1931b}。他劝听众：“所以，爱国的主妇们，明天一早就上街去，到那些到处都在登广告的美妙的大减价里去吧。”\cite{keynes1931b}
% 据: ss_five, ss_vicious, ss_housewives

他还设想了一个极端的情形。大家都不再花钱，把收入全部存起来，结果每个人都会失业，不久也就没有收入可花，“没有一个人会因此多一个便士”\cite{keynes1931b}。五年以后，凯恩斯把这个道理写进了《就业、利息和货币通论》，下面简称《通论》。书的序言署 1935 年 12 月 13 日，书 1936 年在伦敦出版\cite{keynes1936}。照书里的说法，一个人储蓄多少，几乎影响不到自己的收入；可是他的消费是别人的收入，所以所有的人不可能同时储蓄到想要的数目\cite{keynes1936}。“每一次这样靠减少消费来多储蓄的尝试，都会影响收入，使这种尝试必然自我挫败。”\cite{keynes1936}如果人人都想多储蓄，总储蓄为什么可能反而减少？
% 据: ss_stop, ss_penny, gt_preface, gt_title, gt_twosided, gt_defeats

\section{当时的世界}\label{sec:generaltheory-world}

大萧条以前，英国的失业已经高了好几年。\chref{ch:interwar}讲过，1921 年 5 月，英国的行政失业率升到 23.4\%，回到黄金以后也没有降下来。按英格兰银行汇编的数据，1921—1929 年，这个失业率没有一个月低于 8.7\%\cite{boe2017}。萧条来了以后，它在 1932—1933 年最高到了 23.0\%\cite{boe2017}。数据集里另有 Feinstein（1972）的年度估计，以全体平民劳动人口为分母，数字低一些。按这个口径，1932 年有 340 万人失业，占 15.6\%\cite{boe2017}。

\fig{generaltheory_unemp}{英国的失业从 1921 年起一直居高不下，1932 年最重。粗线是每月的行政失业率，1919 年 1—10 月与 1920 年 10 月缺数；琥珀色虚线是 Feinstein（1972）的年度失业率，分母是全体平民劳动人口。圆点标出凯恩斯 1931 年 1 月的广播，竖线标出《通论》出版的 1936 年。数据：英格兰银行《千年宏观数据》（3.1 版）表 M5（Denman 与 Macdonald 1996）与 A50。}{fig:generaltheory_unemp}

1929 年 5 月英国举行大选。前首相劳合·乔治（David Lloyd George，1863—1945）提出了一个开发计划\cite{keynes1929c,hansardpeople32}。凯恩斯与亨德森（Hubert Henderson，1890—1952）合写了一本小册子为它辩护，题为《劳合·乔治做得到吗？》\cite{keynes1929c,loc32}。凯恩斯 1931 年编《劝说集》时，把这篇文字收了进去，题为《扩张的计划》\cite{keynes1929c}。文章开头说，除了回到金本位以前、1924 年那一次短暂的恢复，这个国家有十分之一或更多的劳动人口已经失业了八年，“这在我们的历史上前所未有”\cite{keynes1929c}。按劳工部的统计，参加失业保险的失业者，自 1923 年开始统计以来从没有少于 100 万。1929 年 4 月，失业的工人有 114 万\cite{keynes1929c}。
% 据: pe_ref, pe_title, pe_lg, hh_lloydgeorge, hh_lloydgeorge_pm, loc_henderson, pe_tenth, pe_million, pe_today

失业的代价，他们算得很具体。失业基金每年要付出约 5,000 万镑的现金，1921 年以来已经付出约 5 亿镑，“什么也没有换回来”\cite{keynes1929c}。按 1924 年的生产普查，一个英国工人一年的净产值约 220 镑\cite{keynes1929c}。照这个数算，1921 年以来失业白白浪费的产出约有 20 亿镑\cite{keynes1929c}。劳合·乔治估计，每年 1 亿镑的开发计划能让 50 万人重新就业\cite{keynes1929c}。照他们的计算，这笔钱只相当于 1921 年以来失业损失的 5\%，约合国民收入的 2.5\%\cite{keynes1929c}。
% 据: pe_fund, pe_220, pe_500k, pe_25

财政部的看法正相反。1929 年 4 月 15 日，财政大臣丘吉尔（Winston Churchill，1874—1965）在下院作预算演说\cite{hansard1929,hansardpeople32}。他说，按正统的财政部观点，政府到资金市场上借钱，就成了工业的一个新的竞争者，把本来会被私人企业使用的资源占为己有，还抬高了所有借钱的人要付的利息\cite{hansard1929}。他的政府在这一届议会里花掉或承诺的开发支出已有 4 亿多镑，可是就治疗失业而言，“结果确实令人失望”\cite{hansard1929}。照正统的财政部学说，国家借钱、国家花钱，“事实上，一般来说，只能创造很少的额外就业，不能创造持久的额外就业”\cite{hansard1929}。
% 据: hc29_chancellor, hh_churchill, hc29_view, hc29_400m, hc29_meagre, hc29_doctrine

丘吉尔用划船比赛作比方。冲刺用得好，可以占到好位置，甚至取胜；可是冲刺只是透支全队的体力，失败了，就比一直稳稳地划更糟\cite{hansard1929}。在他看来，劳合·乔治要花 2 亿镑、主要用来修路的计划，不但不能长久地缓解失业，还会严重损害国家的福利\cite{hansard1929}。
% 据: hc29_boat, hc29_spurt, hc29_lg, hc29_injurious

这场争论很快进了一个调查委员会。1929 年 11 月 28 日，金融与工业委员会，人称麦克米伦委员会，第一次听取证词，凯恩斯是委员之一\cite{macmillan1931}。1930 年 4 月 4 日，斯坦普（Josiah Stamp，1880—1941）作证时说，委员会的问卷提出了在“转移”与通胀之间选择的问题。凯恩斯问他这是什么意思。斯坦普答，就是“不管对不对，人们叫作财政部观点的东西”，也就是说，钱从一处拿走，别处就得不到\cite{macmillan1931,loc32}。
% 据: mc_first, mc_keynes, mc_date, mc_alternative, mc_keynesq, mc_stamp, loc_stamp

乘数的想法，是卡恩（Richard Kahn，1905—1989）在 1931 年 6 月的《经济学杂志》上提出的\cite{kahn1931,bnf32}。《通论》说，这个概念是由他第一次引进经济理论的。序言还说，凯恩斯写书时一直依靠他的建议与批评\cite{keynes1936}。卡恩问的是，政府雇人修路，就业一共会增加多少\cite{kahn1931}。修路直接雇的人，他叫作\textbf{初级就业}。这些人花掉新挣的工资，又让生产消费品的人有了活干，他叫作\textbf{次级就业}\cite{kahn1931}。他假定修路的钱靠借款筹集，“银行体系的明智合作”是当然的前提\cite{kahn1931}。
% 据: ka_june, ka_sign, bnf_kahn, gt_kahnfirst, gt_kahnhelp, ka_roads, ka_primary, ka_consumption, ka_secondary, ka_borrowing, ka_banking

\section{有效需求}\label{sec:generaltheory-demand}

《通论》的第一章只有一段。凯恩斯在那里说，古典理论的前提只适用于一种特殊情形，不适用于一般情形；用它去解释经验中的事实，结论会误导人，而且是灾难性的\cite{keynes1936}。序言说，这样得到的是“一种更一般的理论”，人们熟悉的古典理论是它的一个特例\cite{keynes1936}。他说的古典学派，指李嘉图（David Ricardo，1772—1823）的追随者\cite{keynes1936,loc32}。穆勒（John Stuart Mill，1806—1873）与庇古（Arthur Cecil Pigou，1877—1959）都在其中\cite{keynes1936,loc32}。
% 据: gt_special, gt_misleading, gt_general, gt_pigou, loc_ricardo, loc_mill, loc_pigou

就业由什么决定？凯恩斯从企业家的预期讲起。雇用一定数量的工人，产出至少要卖得多少钱，企业家才觉得值得，这是总供给函数；企业家预期实际能卖得多少钱，这是总需求函数\cite{keynes1936}。预期的收入高于所需的收入，企业家就会多雇人。两条函数相交的地方，企业家预期的利润最大，就业量就定在这里。交点上的需求，凯恩斯叫作\textbf{有效需求}\cite{keynes1936}。
% 据: gt_asp, gt_effdemand

照凯恩斯的说法，从萨伊（Jean-Baptiste Say，1767—1832）与李嘉图的时代起，古典经济学家就教导说，“供给创造自己的需求”\cite{keynes1936,loc32}。他把萨伊定律写成，整个产出的总需求价格，在一切产量上都等于它的总供给价格；这等于说，充分就业没有任何障碍\cite{keynes1936}。照 Baumol（1999）的说法，现在一般认为，凯恩斯的这句概括“并没有说得很准确”\cite{baumol1999}。
% 据: gt_say, loc_say, gt_saylaw, bm_keynes

照\chref{ch:say}的区分，凯恩斯的写法更像萨伊恒等式，他真正要质疑的却是均衡版，即工资与利率的调整要多久，会不会把经济带回充分就业。在他看来，有效需求恰好对应充分就业，只是一种特殊情形，要消费倾向与投资诱导正好处在某种关系上才会出现\cite{keynes1936}。一般情形下，两条函数可以在充分就业以下相交。企业家在交点上雇的人，已经是他们预期最有利可图的数目，想做工的人却还有许多找不到工作。
% 据: gt_special28

凯恩斯说，马尔萨斯（Thomas Robert Malthus，1766—1834）曾激烈地反对李嘉图有效需求不可能不足的学说，可是徒劳\cite{keynes1936,loc32}。在他看来，李嘉图征服英国，“就像宗教裁判所征服西班牙一样彻底”\cite{keynes1936}。照凯恩斯的说法，普遍过剩之争（\chref{ch:say}）就此停止，“另一种看法完全消失了”\cite{keynes1936}。
% 据: gt_malthus, loc_malthus, gt_inquisition, gt_ceased

\section{消费函数与乘数}\label{sec:generaltheory-multiplier}

有效需求的一大块是消费。凯恩斯把收入与消费之间的关系叫作消费倾向。他提出一条基本的心理规律：“人们通常而平均地倾向于随着收入的增加而增加消费，但增加得没有收入那么多。”\cite{keynes1936}收入每增加一单位，消费增加的部分，叫作\textbf{边际消费倾向}。
% 据: gt_propensity, gt_psychlaw

把这条规律写成最简单的形式，就是\textbf{凯恩斯交叉}。设消费是收入的线性函数，$C=a+cY$，其中$0<c<1$。$a$是与收入无关的那部分消费，$c$是边际消费倾向。投资$I$暂时看作给定。产出等于计划的支出时，经济处在均衡，
\begin{keyeq}
Y=C+I,\quad C=a+cY\qquad\Longrightarrow\qquad Y^{*}=\frac{a+I}{1-c},\qquad \Delta Y^{*}=\frac{\Delta I}{1-c}.\label{eq:generaltheory-cross}
\end{keyeq}
$1/(1-c)$叫作\textbf{投资乘数}。多出的投资先成为一些人的收入，他们花掉其中的$c$，又成为另一些人的收入，一轮一轮下去，合计是$1+c+c^{2}+\cdots=1/(1-c)$。这是\chref{ch:accounts}的贯穿模型加上的第一组行为方程。图~\ref{fig:generaltheory_cross}(a)画出了它。$45^{\circ}$线上产出等于支出，支出线与它相交的地方就是均衡。

\begin{example}{棉布经济的乘数}\label{ex:generaltheory-cross}
沿用\exref{ex:accounts-thrift}的棉布经济，收入 90，家庭买布 80，投资 10。设消费函数是$C=26+0.6Y$，数字为教学输入。边际消费倾向取 0.6，落在凯恩斯从美国的数字推出的 60\%—70\% 的下端，不过他认为这个数对萧条时期低得不大可信（\secref{sec:generaltheory-expectations}）。（1）均衡收入是多少？（2）乘数是多少？（3）投资从 10 增加到 14，收入与消费各增加多少？
\solution （1）$Y^{*}=(26+10)/(1-0.6)=90$，消费$26+0.6\times90=80$，与\chref{ch:accounts}的账目相同。（2）乘数$1/(1-0.6)=2.5$。（3）$\Delta Y=2.5\times4=10$，收入到 100；消费到$26+0.6\times100=86$，增加 6。一轮一轮地看，多出的 4 先成为收入，花掉 2.4，再花掉 1.44，合计$4+2.4+1.44+\cdots=10$。
\end{example}

卡恩的乘数比这复杂，因为钱在每一轮都会漏出去\cite{kahn1931}。他设每多雇一个人，工资是$W$，利润增加$P$，进口的原料与半成品增加$R$。新就业的人原来领救济$U$，他多出的收入里有$m'$花在本国的消费品上；利润增加的部分里，有$n$花在本国的消费品上\cite{kahn1931}。这样，每雇一个人做初级就业，就带来$k$个人的次级就业，而且一轮接一轮，
\begin{keyeq}
k=\frac{mW+nP}{W+P+R},\qquad m=m'\Bigl(1-\frac{U}{W}\Bigr),\qquad \frac{\text{次级就业}}{\text{初级就业}}=k+k^{2}+\cdots=\frac{k}{1-k}.\label{eq:generaltheory-kahn}
\end{keyeq}
漏出有四处，即进口的原料、进口的成品、救济金的节省，以及利润里没有花掉的部分。卡恩提醒，这是一切都安定下来以后的状态，次级就业要过一段时间才充分显现\cite{kahn1931}。
% 据: ka_wp, ka_r, ka_mwnp, ka_eq, ka_dole, ka_m, ka_settled, ka_lag

\begin{example}{卡恩对英国的估计}\label{ex:generaltheory-kahn}
卡恩估计，英国新增产品的零售价格里，进口的原料与半成品占 1/10，工资占 7/10，利润占 1/5；新就业的人多出的收入，有 1/6 花在进口的成品上，即$m'=5/6$。他认为救济与工资之比$U/W$略低于 1/2，$n$不会低到 1/3\cite{kahn1931}。表~\ref{tab:generaltheory-kahn}的第一行取$U/W=3/7$、$n=3/4$。（1）$m$与$k$各是多少？（2）每 100 人修路，总共能增加多少就业？
\solution （1）$m=\frac{5}{6}\times\bigl(1-\frac{3}{7}\bigr)=\frac{10}{21}$；$k=\frac{10}{21}\times\frac{7}{10}+\frac{3}{4}\times\frac{1}{5}=\frac{1}{3}+\frac{3}{20}=\frac{29}{60}$。（2）次级就业与初级就业之比是$k/(1-k)=29/31\approx0.94$，100 人修路，另有约 94 人就业，合计约 194 人。用凯恩斯交叉的话说，乘数是$1/(1-k)=60/31\approx1.94$，比\exref{ex:generaltheory-cross}的 2.5 小，因为每一轮都有一部分花到了国外，或者只是替换了救济金。
\end{example}

\begin{table}[!tb]\centering\footnotesize
\caption{卡恩对英国的估计：次级就业与初级就业之比}\label{tab:generaltheory-kahn}
\begin{tabularx}{\linewidth}{@{}YYYYY@{}}
\toprule
$U/W$ & $n$ & $m=m'(1-U/W)$ & $k$ & $k/(1-k)$ \\\midrule
3/7 & 3/4 & 10/21 & 29/60 & 29/31 = 0.94 \\
3/7 & 2/3 & 10/21 & 7/15 & 7/8 = 0.88 \\
1/2 & 1/2 & 5/12 & 47/120 & 47/73 = 0.64 \\
1/2 & 1/3 & 5/12 & 43/120 & 43/77 = 0.56 \\
\bottomrule
\end{tabularx}
\datasource{Kahn（1931）第 186 页的表，据扫描页读取；共同的假定是$W/(W+P+R)=7/10$，$P/(W+P+R)=1/5$，$m'=5/6$。最后一列是原表的两位小数。}
\end{table}
% 据: ka_shares, ka_mprime, ka_uw, ka_n, ka_table

卡恩认为，第一行也许偏于乐观；更有把握的是，最后一行偏于保守。他说，若假定实际的比例是 3/4，“也许还是低估了”\cite{kahn1931}。他还说，一个国家越接近封闭，救济金相对于全额工资越少，次级就业相对于初级就业就越多；美国的这个比例比大多数国家大得多\cite{kahn1931}。
% 据: ka_threequarters, ka_table, ka_closed, ka_us

《通论》把乘数与边际消费倾向直接连在一起。凯恩斯在书里举例说，若边际消费倾向是 0.9，乘数就是 10，假定别处的投资没有因此减少，公共工程带来的总就业就是初级就业的 10 倍\cite{keynes1936}。他说得更尖锐。如果财政部把钞票装进旧瓶子，埋进废弃的煤矿，再让私人企业去挖出来，“就不必再有失业”；当然，“盖房子之类的事要明智得多”\cite{keynes1936}。
% 据: gt_kahnfirst, gt_ninetenths, gt_tentimes, gt_noreduction, gt_bottles, gt_nomore, gt_houses

\section{投资、预期与“动物精神”}\label{sec:generaltheory-expectations}

乘数说的是投资变了以后收入怎样变，投资本身又由什么决定？凯恩斯的答案是比较两样东西，一是\textbf{资本的边际效率}，即让一项资本资产预期收益的现值恰好等于它的供给价格的那个贴现率；二是利率\cite{keynes1936}。资本的边际效率高于利率，投资就有利可图。
% 据: gt_mec

资本的边际效率取决于预期，而预期的基础很脆弱。凯恩斯说，对几年以后决定投资收益的那些因素，我们的知识“通常很少，而且往往微不足道”\cite{keynes1936}。在他看来，大多数积极的决定，只能出于\textbf{动物精神}，“一种想要行动而不是不行动的自发冲动”，不是数学期望的结果\cite{keynes1936}。动物精神一旦消沉，自发的乐观动摇，只剩数学期望可以依靠，企业就会萎缩、消亡\cite{keynes1936}。
% 据: gt_precarious, gt_animal, gt_fade

股票市场让事情更复杂。凯恩斯说：“一个国家的资本发展，若成了赌场活动的副产品，这件事多半会做得很糟。”\cite{keynes1936}投资可以剧烈波动，美国的数字就是例子。按库兹涅茨（Simon Kuznets，1901—1985）的估计，1929—1932 年美国的毛资本形成下降了 77.4\%，净资本形成下降了 95.1\%\cite{keynes1936,loc32}。表~\ref{tab:generaltheory-kuznets}是《通论》转引的数字。
% 据: gt_casino, gt_kuznets, loc_kuznets

\begin{table}[!tb]\centering\footnotesize
\caption{美国的资本形成，1925—1933 年（百万美元）}\label{tab:generaltheory-kuznets}
\begin{tabularx}{\linewidth}{@{}p{0.16\linewidth}YYY@{}}
\toprule
年份 & 毛资本形成 & 企业的修理、维护、折旧与损耗 & 净资本形成 \\\midrule
1925 & 30,706 & 7,685 & 23,021 \\
1926 & 33,571 & 8,288 & 25,283 \\
1927 & 31,157 & 8,223 & 22,934 \\
1928 & 33,934 & 8,481 & 25,453 \\
1929 & 34,491 & 9,010 & 25,481 \\
1930 & 27,538 & 8,502 & 19,036 \\
1931 & 18,721 & 7,623 & 11,098 \\
1932 & 7,780 & 6,543 & 1,237 \\
1933 & 14,879 & 8,204 & 6,675 \\
\bottomrule
\end{tabularx}
\datasource{Keynes（1936）第 103 页的表，据扫描页读取。凯恩斯说这些数取自美国国家经济研究局第 52 号公报。毛资本形成已计入企业存货的净变化；净资本形成按库兹涅茨的定义，是前两列之差。}
\end{table}
% 据: gt_kuz_gross, gt_kuz_maint, gt_kuz_net, gt_bulletin

\begin{example}{库兹涅茨的投资数字}\label{ex:generaltheory-kuznets}
用表~\ref{tab:generaltheory-kuznets}。（1）1929—1932 年，毛资本形成与净资本形成各下降了百分之几？（2）1932 年的净资本形成，比 1925—1929 年的平均数低百分之几？（3）凯恩斯说，把年份两两合在一起看，乘数小于 3，大概稳定在 2.5 附近\cite{keynes1936}。按式~\eqref{eq:generaltheory-cross}，这相当于多大的边际消费倾向？
\solution （1）毛资本形成从 34,491 降到 7,780，下降 77.4\%；净资本形成从 25,481 降到 1,237，下降 95.1\%。凯恩斯说，投资据说下降了 75\% 以上，净资本形成下降了 95\% 以上\cite{keynes1936}。（2）1925—1929 年的平均数是 24,434.4，1932 年比它低 94.9\%。（3）$1-1/2.5=0.6$；乘数小于 3，即边际消费倾向小于$1-1/3\approx0.67$。凯恩斯说，这意味着边际消费倾向不超过 60\%—70\%，对繁荣时期说得过去，对萧条时期却低得出奇，在他看来不大可信\cite{keynes1936}。
\end{example}
% 据: gt_overstate, gt_95, gt_pairs, gt_boom

凯恩斯自己并不全信这些数字。他认为，“统计也许多少夸大了投资的下降”，而估计稍有变动，乘数就会差得很多\cite{keynes1936}。
% 据: gt_overstate, gt_substantial

\section{节俭悖论}\label{sec:generaltheory-thrift}

现在可以回到本章的问题了。\exref{ex:accounts-thrift}里，家庭决定少买 10 的布，当年的产出没有变，卖不掉的布记作计划外的存货投资，储蓄等于投资照样成立。恒等式说到这里就停了。凯恩斯交叉接着往下说，存货堆积，企业就会减产。

设家庭在每一个收入水平上都少消费 10，$a$从 26 降到 16。投资仍是 10。按式~\eqref{eq:generaltheory-cross}，新的均衡收入是$(16+10)/0.4=65$，消费是 55，储蓄是$65-55=10$，与原来一样。家庭想多储蓄 10，储蓄却一点也没有增加。收入少了 25，正好是 10 乘以乘数 2.5。这就是\textbf{节俭悖论}。储蓄最后总要等于投资；投资不变，储蓄也就变不了，变的是收入。

投资若也随收入变化，储蓄还会减少。凯恩斯说，个人储蓄等于“决定今天不吃晚饭”，却并不等于决定一周或一年以后买什么\cite{keynes1936}。它压低了今天的生意，却没有刺激为将来准备的生意。对将来消费的预期，大半依据眼下的消费，所以储蓄可能连资本的边际效率也一起压低\cite{keynes1936}。这样，它减少的就不只是眼下的消费需求，还有眼下的投资需求\cite{keynes1936}。
% 据: gt_dinner, gt_depresses, gt_mec210, gt_reduce

庇古不同意这个推论。他打了个比方。决定拔牙，本身并不是决定买假牙，可是很少有人只做前一个决定\cite{pigou1936}。在他看来，一个人储蓄而自己不投资，后果要看银行执行什么政策才能确定\cite{pigou1936}。银行若维持总的货币收入稳定，个人的窖藏就会被别处的支出抵消，资本照样增加\cite{pigou1936}。
% 据: pg_teeth, pg_banking, pg_stable

把凯恩斯的这一点写进模型。设投资是$I=1+0.1Y$，收入为 90 时投资仍是 10。家庭同样少消费 10，收入降到约 56.7，投资与储蓄都降到约 6.7（\exref{ex:generaltheory-thrift}）。图~\ref{fig:generaltheory_cross}(b)画出了两种情形。人人都想多储蓄，总储蓄反而减少了。

\fig{generaltheory_cross}{凯恩斯交叉与节俭悖论。(a) 支出线$C+I$与$45^{\circ}$线相交于收入 90；投资从 10 增加到 14，支出线上移 4，收入增加 10，乘数是 2.5。(b) 储蓄线$S=Y-C$与投资线相交的地方是均衡。家庭在每个收入上都想多储蓄 10，储蓄线上移。投资固定在 10 时，收入从 90 降到 65，储蓄仍是 10；投资随收入变化（$I=1+0.1Y$，琥珀色虚线）时，收入降到约 56.7，储蓄降到约 6.7。数字为教学输入，$C=26+0.6Y$。}{fig:generaltheory_cross}

\begin{example}{节俭的结局}\label{ex:generaltheory-thrift}
接\exref{ex:generaltheory-cross}，消费函数从$C=26+0.6Y$变成$C=16+0.6Y$。（1）投资固定在 10，新的收入、消费与储蓄各是多少？（2）投资是$I=1+0.1Y$，原来的收入是多少？新的收入、投资与储蓄各是多少？（3）\exref{ex:accounts-thrift}里第一年的储蓄是 20，这与（1）（2）的结果矛盾吗？
\solution （1）$Y=(16+10)/(1-0.6)=65$，$C=16+0.6\times65=55$，$S=65-55=10$。（2）原来$Y=(26+1)/(1-0.6-0.1)=90$。新的$Y=(16+1)/0.3\approx56.7$，$I=1+0.1\times56.7\approx6.7$，$C=16+0.6\times56.7=50$，$S=56.7-50\approx6.7$。（3）不矛盾。第一年的 20 里有 10 是卖不掉的布。企业减产以后，计划外的存货消失了，储蓄回到投资的水平。投资固定时是 10，投资随收入下降时只有约 6.7。
\end{example}

凯恩斯在书里说得更重：“我们越有美德，越是决心节俭，在国家与个人的财政上越是顽固地守旧，利率相对于资本的边际效率上升时，我们的收入就越要下降。”\cite{keynes1936}照斯密（Adam Smith，1723—1790）的说法，每年省下的东西照常被消费掉，只是消费的人不同（\chref{ch:smith}）\cite{smith1776,loc32}。省下的钱要真被用来雇人、添置资本，这句话才成立。凯恩斯的回答是，一个人节俭是美德，大家同时节俭，需求就可能不足。决定储蓄有多少的，最后是投资与收入，不是节俭的意愿。
% 据: gt_virtuous, loc_smith

\begin{compare}{两种储蓄—投资}{古典的储蓄—投资}{凯恩斯的储蓄—投资}
\cmprow{储蓄等于投资}{古典学派也承认二者相等，多数人还把这一点推得太远，以为个人每多储蓄一笔，必然带来相应的一笔投资\cite{keynes1936}}{必然相等，因为两者都等于收入超过消费的部分\cite{keynes1936}}
\cmprow{谁来调整}{利率。储蓄是可投资资源的供给，投资是需求，利率是使二者相等的价格\cite{keynes1936}}{收入。想多储蓄，收入就下降，直到储蓄回到投资的水平}
\cmprow{多储蓄的结果}{利率下降，投资增加，资本积累更快}{收入与就业下降；投资随收入或预期下降时，储蓄反而减少}
\cmprow{利率由什么决定}{储蓄与投资}{货币数量与流动性偏好（\secref{sec:generaltheory-liquidity}）}
\cmpsum{两者都说储蓄等于投资。古典理论认为，利率让储蓄的供给与投资的需求相等。凯恩斯认为二者总是相等，调整的是收入，收入一直变到人们想储蓄的数目恰好等于投资为止。}
\end{compare}
% 据: gt_identity63, gt_toofar, gt_classical175

\begin{pitfall}[节俭悖论说储蓄总是坏事]
凯恩斯 1931 年的广播并没有这样说。他说，多储蓄“在某些情况下”完全对\cite{keynes1931b}。储蓄的目的，是腾出劳动去生产房屋、工厂、道路、机器之类的资本品；可是已有大批失业者可供这些用途的时候，储蓄只是增加了失业\cite{keynes1931b}。在《通论》里，有效需求对应充分就业是一种特殊情形，那时多储蓄确实可以变成多投资。悖论要成立，要有两个条件，一是大家同时多储蓄，二是投资不因此增加。
\end{pitfall}
% 据: ss_wrong, ss_object, ss_surplus, gt_special28

\section{流动性偏好与利率}\label{sec:generaltheory-liquidity}

古典理论说，利率让储蓄与投资相等。在凯恩斯看来，让储蓄等于投资的是收入，利率由什么决定，就要另作解释。他的回答是：“利率是在一段特定的时间里放弃流动性的报酬。”\cite{keynes1936}它不是让投资需求与放弃当前消费的意愿达成均衡的“价格”\cite{keynes1936}。
% 据: gt_reward, gt_notprice

人们为什么要持有不生息的货币？凯恩斯举了三个动机，即交易、预防与投机\cite{keynes1936}。前两者的需求主要随收入变化。投机的需求，取决于现行利率与人们对将来利率的预期\cite{keynes1936}。他把货币$M$分成两部分，$M_1$满足前两个动机，$M_2$满足投机的动机，
\begin{keyeq}
M=M_1+M_2=L_1(Y)+L_2(r).\label{eq:generaltheory-liquidity}
\end{keyeq}
$L_1$主要取决于收入$Y$，$L_2$主要取决于利率$r$与预期的关系。照凯恩斯的说法，流动性偏好得以存在的必要条件，是对将来的利率不确定\cite{keynes1936}。持有债券的人担心利率上升、债券跌价，就宁愿持有货币。\chref{ch:say}里，穆勒说过人们在某些时候“宁愿持有货币，而不愿持有任何别的商品”\cite{jsmill1844}；凯恩斯把这种偏好写成了一个随利率变化的函数。
% 据: gt_motives, gt_speculative, gt_l1l2, gt_uncertainty, jm_possess

凯恩斯还提到一种极端情形。利率降到一定水平以后，几乎人人都宁愿持有现金，流动性偏好可能变成“实际上绝对的”，货币当局就控制不了利率了\cite{keynes1936}。他又说，这种情形“我至今不知道有过例子”\cite{keynes1936}。后人把它叫作流动性陷阱，用一条曲线来画（\chref{ch:islm}）。
% 据: gt_absolute, gt_noexample

这个利率理论很快招来了反驳。1937 年 6 月，俄林（Bertil Ohlin，1899—1979）在《经济学杂志》上说，利率“只是信用的价格”，由信用的供给与需求决定\cite{ohlin1937b,loc32}。这种看法通常称为\textbf{可贷资金理论}。同一期上，凯恩斯答复说，按俄林的定义，信用的净供给与储蓄的数量“完全是同一个东西”，结论又回到了古典学说\cite{keynes1937}。
% 据: oh2_june, oh2_price, loc_ohlin, k37_title, k37_sign, k37_same

\begin{compare}{可贷资金与流动性偏好}{可贷资金（俄林，1937）}{流动性偏好（凯恩斯，1936）}
\cmprow{利率是什么}{信用的价格\cite{ohlin1937b}}{放弃流动性的报酬\cite{keynes1936}}
\cmprow{由什么决定}{信用的供给与需求，银行体系也能影响}{货币数量与持有货币的意愿，即$L_1(Y)+L_2(r)$}
\cmprow{多储蓄会怎样}{信用的供给增加，利率下降}{收入先下降；交易需要的货币少了，利率才可能随之下降}
\cmprow{增加货币会怎样}{信用增加，利率下降}{利率可能下降，但流动性偏好可能吸收掉增加的货币}
\cmpsum{两者都承认银行与货币能影响利率。分歧在于储蓄的意愿是不是利率的直接决定因素。可贷资金理论说是，凯恩斯说利率决定于存量的货币与持有它的意愿。}
\end{compare}

增加货币，能不能在有失业时增加产出？这是\chref{ch:mercantilism}里约翰·劳（John Law，1671—1729）1705 年的主张\cite{law1705,loc32}。照 Velde（2003）的说法，劳的名声在 1930 年代上升了\cite{velde2003}。凯恩斯的回答是，可能，但中间有几道关口。他提醒，从杯子到嘴边，“可能会有好几次失手”\cite{keynes1936}。货币增加了，利率却未必下降，因为人们可能更想持有货币。利率降了，投资却未必增加，因为资本的边际效率可能降得更快。投资增加了，就业也未必增加，因为消费倾向可能在下降\cite{keynes1936}。
% 据: loc_law, velde_1930s, gt_slips, gt_liqfalls, gt_mecfalls, gt_propensityfalls

关口都过了，增加的有效需求也会一部分变成产出与就业，一部分变成物价。凯恩斯把数量论改写成：“只要还有失业，就业就与货币数量同比例变化；到了充分就业，物价就与货币数量同比例变化。”\cite{keynes1936}他马上说明，这只是加了许多简化假定以后的说法\cite{keynes1936}。凯恩斯把普通劳动雇用一小时算作一个劳动单位，一个劳动单位的货币工资叫作\textbf{工资单位}\cite{keynes1936}。实际上，工资单位在充分就业以前就会上涨，物价随着就业的增加逐步上升\cite{keynes1936}。
% 据: gt_partly, gt_qt, gt_simplify, gt_labourunit, gt_wageunitdef, gt_wageunit

\section{同时代的发现}\label{sec:generaltheory-parallel}

有效需求不只在剑桥被发现。瑞典的经济学家走的是自己的路，卡莱茨基在波兰也算出了相近的结果。它们与《通论》有什么相同，有什么不同？

\begin{case}{瑞典 · 1931—1937 · 斯德哥尔摩学派}
1937 年 3 月，俄林在《经济学杂志》上向英语读者介绍瑞典经济学家的工作\cite{ohlin1937}。他说，萧条刚开始，瑞典的经济学家就用一套不同于教科书价格理论的工具，研究就业、产出与物价的变动；它与《通论》是对同一组问题的“两次独立的进攻”\cite{ohlin1937}。他把这条路追溯到维克塞尔（Knut Wicksell，1851—1926），说维克塞尔已经与“供给创造自己的需求”的萨伊学说决裂（\chref{ch:wicksell}）\cite{ohlin1937,loc32}。
% 据: oh_title, oh_march, oh_early, oh_two, oh_say, loc_wicksell

照俄林的说法，1931 年底，瑞典的失业委员会请四位经济学家写专题报告，问的是在大量资源闲置时，各种政策有什么效果\cite{ohlin1937}。四份报告在 1933 年底到 1934 年初出版，其中有缪尔达尔（Gunnar Myrdal，1898—1987）的《公共财政政策的经济效果》与俄林自己的书\cite{ohlin1937,loc32}。照 Berg 与 Jonung（1998）的说法，缪尔达尔还为财政大臣维格福斯（Ernst Wigforss，1881—1977）1933 年的预算写了一份报告，给逆周期的财政政策提供了理论基础\cite{berg1998,loc32}。
% 据: oh_1931, oh_unused, oh_published, oh_myrdal, loc_myrdal, bj_myrdal, loc_wigforss

两边也有差别。照俄林的说法，委员会 1935 年的总结报告与《通论》讨论的是同一片领域；它很少讨论利率怎样决定，却详细分析了凯恩斯几乎完全没有理会的“摩擦性”失业\cite{ohlin1937}。这些人的看法相当一致，俄林认为可以称为“斯德哥尔摩学派”\cite{ohlin1937}。在利率上，他后来走向了可贷资金理论（\secref{sec:generaltheory-liquidity}）。
% 据: oh_1935, oh_field, oh_frictional, oh_school
\end{case}

卡莱茨基（Michał Kalecki，1899—1970）的论文，1933 年 10 月在莱顿的计量经济学会会议上宣读，1935 年刊在《计量经济学》上\cite{kalecki1935,bnf32}。他把资本家的实际毛利润分成消费的与积累的两部分，资本家的消费又分成固定的部分与随利润变化的部分\cite{kalecki1935}。于是，毛利润与固定消费加积累之和成正比\cite{kalecki1935}。第一个等式就是\chref{ch:marx}的式~\eqref{eq:marx-kalecki}，即毛利润等于毛投资加资本家的消费。比例系数在形式上与乘数相同。
% 据: kc_toc, kc_leyden, kc_sign, bnf_kalecki, kc_parts, kc_elastic, kc_proportionate

卡莱茨基也用钱的循环说明支出怎样变成收入。一些资本家从储蓄存款里取出一笔钱，或者向中央银行借来，去建新的设备\cite{kalecki1935}。同一年里，这笔钱以利润的形式到了别的资本家手里，又存回银行，“循环就这样闭合了”\cite{kalecki1935}。他还声明，他假定存在“一支失业的后备军”\cite{kalecki1935}。与凯恩斯一样，这条路要靠闲置的资源。
% 据: kc_withdraw, kc_circle, kc_reserve

政策走在理论的前面。\chref{ch:alternatives}讲过，瑞典、日本与德国的扩张政策，都在《通论》出版以前开始。照 Berg 与 Jonung 的说法，瑞典 1931—1932 年的货币计划产生在一个“前凯恩斯的世界”里\cite{berg1998}。照 Shizume（2009）的记述，有些观察者称赞高桥是清（1854—1936）是凯恩斯经济学的一位成功的先驱\cite{shizume2009,loc32}。有效需求理论后来给这些做法提供了一种解释。资源闲置时，赤字支出与中央银行的信用，可以通过乘数增加收入与就业，不只推高物价。
% 据: bj_prekeynes, sz_pioneer, loc_takahashi

\section{两场争论}\label{sec:generaltheory-debate}

\subsubsection{1929 年的财政部观点错了吗}
分歧在于，在有大量失业的时候，国家借钱修路，能不能增加总就业，还是只把就业从私人企业挪到公共工程。

财政部一方最有力的论证是资金的竞争。丘吉尔说，政府借钱要与工业争夺资源，还会抬高利率\cite{hansard1929}。他拿自己政府的经验作证据，4 亿多镑的开发支出，对失业的效果令人失望\cite{hansard1929}。他还援引一部经济史的标准著作说，1845 年前后的铁路热里，人们借了大笔的钱赶修铁路，接着就是萧条与金融危机\cite{hansard1929}。
% 据: hc29_view, hc29_400m, hc29_meagre, hc29_railway, hc29_crisis

凯恩斯与亨德森的回答有三层。第一，这个论证若成立，也同样适用于私人企业新建工厂，那么任何新的投资都不能增加就业\cite{keynes1929c}。第二，开发支出所需的资金有三个来源，即现在用来养活失业者的储蓄，因为信用不足而白白浪费的储蓄，以及减少的对外贷款\cite{keynes1929c}。第三，在大家都已经就业以后还要扩张，才会有通胀\cite{keynes1929c}。眼下先要吸收“大量的通缩性闲置”\cite{keynes1929c}。
% 据: pe_morris, pe_three, pe_inflation, pe_slack

他们也作了让步。英格兰银行若执行紧缩政策，阻止银行信用的任何扩张，就能打乱最周密的计划，让财政部出钱的支出确实以别的企业为代价\cite{keynes1929c}。所以，财政部观点对不对，取决于中央银行怎样做。卡恩 1931 年的估计，正是在“银行体系的明智合作”的前提下做出的\cite{kahn1931}。
% 据: pe_bank, pe_bestlaid, ka_banking

到了《通论》，凯恩斯自己也列出了几种会抵消公共工程效果的因素\cite{keynes1936}。筹资的方式，以及就业增加、物价上涨以后多需要的现金，可能抬高利率，减少别处的投资，除非货币当局采取相反的措施\cite{keynes1936}。在常见的混乱心理下，政府的计划还可能通过对“信心”的影响，提高流动性偏好，或者压低资本的边际效率，同样会减少别处的投资，除非另有措施加以抵消\cite{keynes1936}。这个问题后来成了财政政策争论的老题目（\chref{ch:budget}）。
% 据: gt_offset_rate, gt_confidence

\subsubsection{一般理论，还是工资刚性的经济学}
另一场争论围绕货币工资。分歧在于，削减货币工资能不能增加就业。《通论》的书名说它是一般理论；批评者说，它的结论靠的是货币工资不肯下降。

凯恩斯一方的论证，在书中讨论货币工资变动的一章。他认为，削减货币工资不一定增加就业。工资降了，收入与物价随着下降，交易需要的现金就少了，在其他条件不变时利率会下降，从而有利于投资\cite{keynes1936}。这样，削减工资对利率的作用，与不动工资、增加货币数量相同；作为实现充分就业的办法，它也受到同样的限制\cite{keynes1936}。他说，没有理由相信灵活的工资政策能够维持持续的充分就业，正如没有理由相信公开市场的货币政策单靠自己就能做到\cite{keynes1936}。
% 据: gt_wagecash, gt_wagerate, gt_samelimits, gt_flexible, gt_unaided

在他看来，灵活的工资政策与灵活的货币政策，在分析上是一回事，都是改变以工资单位计算的货币数量\cite{keynes1936}。差别在实行上。除非工资由法令决定，没有办法让各类劳动的工资一齐下降，结果只能是一连串零碎的变动，很可能要经过浪费而惨痛的斗争才完成，议价地位最弱的人吃亏最多\cite{keynes1936}。改变货币数量，却已在大多数政府的权力之内，靠公开市场政策就能做到\cite{keynes1936}。所以他说，考虑到人性与我们的制度，宁要灵活的工资政策、不要灵活的货币政策的，“只能是愚蠢的人”，除非他能指出前者有后者得不到的好处\cite{keynes1936}。他还说，降低工资单位会按比例加重债务的负担，增加货币数量的作用正相反\cite{keynes1936}。
% 据: gt_samething, gt_uniform, gt_struggles, gt_openmarket, gt_foolish, gt_debt

庇古的回答在 1936 年 5 月的《经济学刊》上\cite{pigou1936}。他先说，爱因斯坦为物理学做到的，正是凯恩斯以为自己为经济学做到的事；可是爱因斯坦没有暗示牛顿是一伙无能的人\cite{pigou1936}。照他的读法，凯恩斯的结论与通常的看法没有什么不同，也承认长期里只要货币工资能自由变动，经济会趋向充分就业\cite{pigou1936}。他也让了一步。像英国 1925—1929 年那样失业严重的时候，用银行政策推高物价，有很强的理由比削减货币工资更好。他说，理由是人们熟悉的，凯恩斯在书里也提到过。在一个民主国家，很难普遍地大幅削减货币工资；即使做到了，私人与国家的债务的实际负担也会加重\cite{pigou1936}。
% 据: pg_title, pg_may, pg_einstein, pg_bunglers, pg_conclusion, pg_longperiod, pg_case, pg_strong, pg_democracy, pg_debts

分歧落在过程上。庇古认为，只要银行政策与实际推行过的相像，削减货币工资“必然降低实际工资、增加就业”\cite{pigou1936}。他还说，1925—1929 年削减货币工资的理由，比后来萧条的几年更充分\cite{pigou1936}。在他看来，凯恩斯是一位艺术家，“向月亮射箭”，箭法如何另当别论，技艺却人人都能欣赏\cite{pigou1936}。后来的争论问的是，工资与物价若能完全灵活，经济会不会自己回到充分就业（第~\ref{ch:islm}、\ref{ch:naturalrate}~章）。
% 据: pg_must, pg_stronger, pg_arrows

\section{数据核对与反例}\label{sec:generaltheory-check}

如果萧条来自有效需求不足，产出下降时，下降的应当主要是投资、出口这类不随国内收入变化的支出，消费只随收入变化。英国的数据有合有不合。按英格兰银行数据集的实际值，1929—1932 年英国的实际 GDP 下降了 5.3\%\cite{boe2017}。实际投资下降了 14.1\%，出口量下降了 32.2\%\cite{boe2017}。实际消费反而增加了 2.1\%\cite{boe2017}。

\fig{generaltheory_uk}{英国的萧条与复苏，1929—1938 年。(a) 实际 GDP、消费、投资与出口量，1929 年$=100$。1932 年出口量最低，消费一直没有低于 1929 年；1932 年以后投资领先回升。(b) 英格兰银行的贴现率，1932 年 6 月 30 日降到 2\%，一直保持到 1939 年 8 月。数据：英格兰银行《千年宏观数据》（3.1 版）表 A1 与 D1。}{fig:generaltheory_uk}

消费的上升，简单的消费函数解释不了，它只说消费随收入同向变化。出口的下跌在凯恩斯交叉里与投资的下跌作用相同，都是不随国内收入变化的支出。出口量跌得最多，说明英国的萧条与世界经济连在一起（\chref{ch:depressionworld}）。美国不同，投资跌得最厉害（\tabref{tab:generaltheory-kuznets}）。照卡恩的说法，美国的次级就业比例比大多数国家大得多\cite{kahn1931}。他对英国的估计，次级就业是初级就业的 0.56 到 0.94\cite{kahn1931}。凯恩斯对美国估计的乘数是 2.5 左右\cite{keynes1936}。
% 据: ka_us, ka_table, gt_pairs

复苏来得不算慢。1932 年 6 月 30 日，英格兰银行把贴现率降到 2\%，一直保持到 1939 年 8 月\cite{boe2017}。1932—1937 年，英国的实际投资增加了 47.7\%，实际 GDP 增加了 22.9\%\cite{boe2017}。按 Feinstein 的口径，失业率从 1932 年的 15.6\% 降到 1937 年的 7.8\%\cite{boe2017}。1938 年它又升到 9.3\%。同一年美国的衰退，是\chref{ch:islm}的题目。

复苏也没有到达每一个地方。1936 年 11 月 4 日，贾罗选出的议员威尔金森（Ellen Wilkinson，1891—1947）向下院递交了贾罗居民的请愿书\cite{hansard1936,hansardpeople32}。请愿书说，过去 15 年，贾罗经历了“全镇历史上从未有过的工业萧条”，“它的造船厂关了门”\cite{hansard1936}。从前雇用 8,000 人的地方，现在只有 100 人在一个临时项目上做工\cite{hansard1936}。
% 据: hc36_ref, hc36_petition, hh_wilkinson, hh_wilkinson_jarrow, hc36b_member, hc36_15, hc36_shipyard, hc36_8000

同一天，下院还就贾罗提出了两个紧急质询。威尔金森问首相，收到了多少份关于贾罗处境的决议\cite{hansard1936}。另一位议员问贸易部大臣，知不知道贾罗有大批失业的造船工人\cite{hansard1936}。贸易部大臣朗西曼（Walter Runciman，1870—1949）答复了这个质询\cite{hansard1936,hansardpeople32}。他说，按劳工部的最新资料，9 月 28 日贾罗有 1,185 名失业的造船工人，另有 818 名普通劳工失业\cite{hansard1936}。他又转述首相 10 月 29 日给威尔金森的答复，说贾罗的失业状况“虽然还远不能令人满意，最近几个月已有改善”\cite{hansard1936}。
% 据: hc36b_ref, hc36b_pm, hc36b_magnay, hc36b_runciman, hh_runciman, hh_runciman_bt, hc36b_1185, hc36b_member, hc36b_improved

\chref{ch:intro}的主线说，大萧条让萨伊定律失去了说服力。英国的数据说明了原因。失业率连续十几年居高不下，工资与利率的调整没有把经济带回充分就业。是调整太慢，还是调整本身就不起作用，正是上一节的争论。

\begin{deeper}[史料说明]
\textbf{《通论》的版本。}正文据 archive.org 上印度数字图书馆所藏 1936 年麦克米伦公司初版的识别文本，页码照原书。第 25—26 页的总需求与总供给、第 103 页的表、第 167 页利率的定义、第 199 页的式子与第 296 页数量论的改写，识别文本有错乱，按扫描页读取。书的定价五先令，见卷首凯恩斯著作的书目。出版的具体日子，本书没有找到原始的出处，正文只写 1936 年。

\textbf{凯恩斯的两篇短文。}《储蓄与花费》与《扩张的计划》都据 1931 年《劝说集》的识别文本。前者注明 1931 年 1 月广播，讲稿登在 1 月 14 日的《听众》上。后者注明出自与亨德森合写的小册子（1929 年 5 月），小册子本身没有取得。识别文本丢掉了镑的符号，钱数按扫描页核对。

\textbf{两处不一致。}凯恩斯与亨德森引丘吉尔的预算演说，作“正统的财政部教条”，《议会议事录》印的是“正统的财政部学说”。丘吉尔说劳合·乔治要花 2 亿镑，凯恩斯与亨德森说的是每年 1 亿镑，两者的口径没有说明。议事录里的劳合·乔治，是“卡那封自治市选出的议员”，他的人物页写着这个选区。丘吉尔的名字只出现在演说的第一页（第 27 栏），引文在第 53—56 栏。1929 年的白皮书（Cmd. 3331）没有取得。

\textbf{数字的来源与核对。}卡恩的分数与式子据扫描页读取。库兹涅茨的数字只读到《通论》的转引，美国国家经济研究局第 52 号公报没有取得。凯恩斯说 1932 年的净资本形成比 1925—1929 年的平均数低“不少于 95\%”，按表算是 94.9\%。

\textbf{失业的口径。}凯恩斯与亨德森说的不少于 100 万，是劳工部统计的参加失业保险的失业者；1929 年 4 月的 114 万，原文没有说明口径。Feinstein 的年度估计以全体平民劳动人口为分母，1929 年的失业人数约 150 万\cite{boe2017}。表 M5 的月度行政失业率，数据集没有说明它覆盖哪些工人，Denman 与 Macdonald（1996）的原文本书没有取得。

\textbf{转述与间接的证词。}缪尔达尔 1933 年给预算写的报告，据 Berg 与 Jonung 的转述；斯德哥尔摩学派的四份报告与 1935 年的总结报告，据俄林 1937 年的介绍，都没有读到原书。称高桥是清为凯恩斯先驱的观察者，据 Shizume 的转述。庇古说凯恩斯允许长期里的充分就业，是庇古自己的读法。贾罗的请愿书是居民的话，由威尔金森递交；朗西曼的数字来自劳工部。

\textbf{生卒年。}凯恩斯、亨德森、斯坦普、李嘉图、穆勒、庇古、萨伊、马尔萨斯、库兹涅茨、约翰·劳、维克塞尔、俄林、缪尔达尔、维格福斯与高桥是清据美国国会图书馆的规范档；卡恩与卡莱茨基据法国国家图书馆的规范档；劳合·乔治、丘吉尔、威尔金森与朗西曼据英国议会历史议事录的人物页。
\end{deeper}
% 据: gt_title, gt_copyright, gt_price, ep_copyright, ss_listener, pe_ref, pe_500k, pe_dogma, hc29_doctrine, hc29_lg, hh_lloydgeorge, hh_lloydgeorge_seat, hc29_ref, hc29_chancellor, ka_eq, ka_table, gt_bulletin, gt_95, pe_million, pe_today

\section{今天还剩什么}\label{sec:generaltheory-legacy}

\textbf{从总支出出发看经济。}产出与就业由有效需求决定，这个问题的提法留了下来。\chref{ch:accounts}的国民核算按支出法把 GDP 分成消费、投资、政府支出与净出口，本章的交叉只是其中最简单的一种行为读法。消费函数后来被永久收入与生命周期的理论改写（\chref{ch:consumption}）。

\textbf{乘数与财政政策。}乘数有多大，至今仍是财政政策争论的中心（\chref{ch:multipliers}）。卡恩已经看到，开放经济、救济制度与银行的态度都会改变它。战时，同一套总需求的算法被反过来用于抑制需求（\chref{ch:warfinance}）。

\textbf{教科书的凯恩斯。}后来的教科书把《通论》写成几条方程与曲线（\chref{ch:islm}），又与价格调整的理论合在一起（第~\ref{ch:phillips}、\ref{ch:naturalrate}~章）。下面的对照框只比较本章的交叉与凯恩斯本人，曲线留给下一章。

\begin{compare}{凯恩斯本人与凯恩斯交叉}{凯恩斯本人（1936）}{凯恩斯交叉（本章的教学模型）}
\cmprow{投资}{取决于资本的边际效率，也就是对将来收益的预期与动物精神，可以剧烈波动}{一个给定的数$I$，或者简单地随收入变化}
\cmprow{货币工资}{削减工资不一定增加就业，灵活的工资政策维持不了充分就业}{价格与工资不变，不讨论}
\cmprow{利率与货币}{流动性偏好决定利率，对将来的利率不确定是它存在的条件}{没有货币，也没有利率}
\cmprow{时间}{卡恩说次级就业要过一段时间才充分显现}{均衡一步到位}
\cmpsum{交叉是凯恩斯理论里最容易计算的一块。它省掉的预期、货币与工资，正是后面各章要补回来的东西。}
\end{compare}

\begin{takeaways}
\begin{itemize}
\item 1921—1938 年，英国的行政失业率没有一个月低于 8.7\%。1929 年，财政部认为国家借钱只能创造很少的额外就业；凯恩斯与亨德森认为，有大量闲置时，开发支出可以增加总就业。
\item 有效需求是总需求函数与总供给函数的交点。它可以对应充分就业以下的就业量，充分就业只是一种特殊情形。
\item 按式~\eqref{eq:generaltheory-cross}，投资变动$\Delta I$，收入变动$\Delta I/(1-c)$。卡恩的式~\eqref{eq:generaltheory-kahn}加上了进口、救济与利润里存起来的部分，他对英国的估计，次级就业是初级就业的 0.56 到 0.94。
\item 节俭悖论说的是，人人想多储蓄，投资不变时储蓄不变而收入下降；投资随收入或预期下降时，储蓄反而减少。
\item 按式~\eqref{eq:generaltheory-liquidity}，利率是放弃流动性的报酬，由货币数量与流动性偏好决定。增加货币能否增加就业，要过利率、投资与消费倾向三道关口。
\item 凯恩斯认为灵活的工资政策维持不了充分就业，庇古认为在现实的银行政策下削减货币工资会增加就业；斯德哥尔摩学派与卡莱茨基独立地走到了相近的结论。
\end{itemize}
\end{takeaways}

\begin{selfcheck}
\begin{enumerate}
\item 在\exref{ex:generaltheory-cross}里，边际消费倾向从 0.6 提高到 0.75，$a$不变。均衡收入与乘数各是多少？
\item 用表~\ref{tab:generaltheory-kahn}的最后一行，验算$k=43/120$与次级就业的比例。
\item 为什么在\exref{ex:generaltheory-thrift}的（1）里，家庭少消费 10，收入却下降了 25？
\item 凯恩斯说灵活的工资政策与灵活的货币政策在分析上是一回事。这是什么意思？他为什么还是偏向后者？
\item 1929—1932 年，英国的实际消费增加了 2.1\%，实际 GDP 下降了 5.3\%。这与简单的消费函数相合吗？
\item 1929 年财政部与凯恩斯争论的是什么？凯恩斯与亨德森承认，在什么条件下财政部的说法会成立？
\end{enumerate}
\end{selfcheck}

\begin{answers}
\item $Y^{*}=(26+10)/(1-0.75)=144$，乘数$1/(1-0.75)=4$。
\item $m=\frac{5}{6}\times\frac{1}{2}=\frac{5}{12}$。$k=\frac{5}{12}\times\frac{7}{10}+\frac{1}{3}\times\frac{1}{5}=\frac{35}{120}+\frac{8}{120}=\frac{43}{120}$，$k/(1-k)=43/77\approx0.56$。
\item 收入下降，消费也随之减少，每减少一轮收入，又少花其中的 0.6。合计是$10\times2.5=25$。
\item 两者都改变以工资单位计算的货币数量，所以对利率与就业的作用相同，受到的限制也相同。差别在实行上。除非工资由法令决定，各类劳动的工资没法一齐下降，只能经过零碎的变动与惨痛的斗争，议价地位最弱的人吃亏最多；改变货币数量，却在大多数政府的权力之内。降低工资单位还会加重债务的负担。庇古也承认这些理由，他概括的常见理由有两条，一是在民主国家很难普遍削减工资，二是削减了又会加重债务的实际负担。
\item 不相合。简单的消费函数说消费随收入同向变化；英国的消费在收入下降时上升了，要用别的原因解释。
\item 争论国家借钱的开发支出能不能增加总就业。英格兰银行若执行紧缩政策，阻止银行信用的扩张，财政部出钱的支出就会以别的企业为代价。
\end{answers}

1931 年 1 月，凯恩斯在广播里说，每储蓄五先令，就让一个人失业一天。五年以后出版的《通论》，定价也是五先令\cite{keynes1936}。序言说，难处不在新的想法，“而在于摆脱旧的想法”\cite{keynes1936}。书里有消费函数与乘数，也有动物精神、流动性偏好与货币工资。后来，这本书被写成了几条方程与一张图。一张图装下了多少凯恩斯，又丢掉了什么？
% 据: ss_five, gt_price, gt_difficulty
